\documentclass[10pt,a4paper]{amsart}
\usepackage[margin=19mm]{geometry}
\usepackage{amsmath,amssymb,mathtools}
\usepackage[scr=rsfs,cal=euler]{mathalfa}
\usepackage{xfrac}
\usepackage[unicode,hidelinks,bookmarks=false]{hyperref}
\newcommand{\ie}{i.e.\ }
\newcommand{\cf}{cf.\ }
\newcommand{\resp}{respectively\ }
\newcommand{\Subsection}{\subsection}
\providecommand{\nobreakdash}{\nobreak\hbox{-}}
\setlength{\emergencystretch}{2em}
\multlinegap0pt
\usepackage{amssymb}
\usepackage[shortlabels]{enumitem}
\newtheorem{proposition}{Proposition}
\newcommand{\tr}{\operatorname{tr}}
\title{The $(n-2,2)$-Spectrum of a Graph}
\author{Boris Shapiro}
\date{Source: arXiv:2605.17501v2 (CC BY 4.0)}
\begin{document}
\maketitle

\noindent\emph{Source and licence.} The $(n-2,2)$-Spectrum of a Graph. Version arXiv:2605.17501v2 (CC BY 4.0), distributed by arXiv under the Creative Commons Attribution 4.0 International licence, http://creativecommons.org/licenses/by/4.0/, as recorded in the arXiv licence metadata for this exact version and shown on the abstract page https://arxiv.org/abs/2605.17501v2. Full licence notice: \texttt{licenses/arxiv-2605.17501-CC-BY-4.0.txt}.\par\smallskip

\noindent\textit{Editorial scope.} Counted unit: the article's proposition proposition (source0.tex lines 793--822). The article's complete original proof of it is delivered with it (source0.tex lines 824--860, one \texttt{proof} environment of 37 lines). No mathematics is written, repaired or re-derived: the statement and the proof are the article's own lines, in the article's own notation, and no retained line is substituted or re-flowed.  No further source-local context is needed: every cross-reference of every retained line resolves inside the two delivered spans.  Nothing was omitted from any retained span, so the package carries no omission record.  No retained line contains a citation, so the wrapper carries no bibliography and no bibliography entry.  No retained line contains a figure environment, an image-inclusion command or drawing code, so no image is delivered and the package declares no asset dependency.  Editorial formatting only, in addition: an AMS \texttt{amsart} preamble that restates the article's own declarations and macros used by the retained lines, namely the environments \texttt{proposition} on separate counters, together with the standard packages those lines need.  The retained environments print in this document as Proposition 1 in order of appearance, whereas the article prints the same items with the numbering of its own full document.  The complete native source of the article as distributed in the arXiv e-print for this exact version is bundled unchanged as \texttt{sources/200800/source0.tex}.
\begin{proposition}[Reduction modulo Laplacian data]
Let
\[
\ell_j(G)=\tr(L_G^j),\qquad
p(G)=\frac{\ell_2(G)-4m}{2},
\]
and put
\[
a(G):=s(G)-t(G)
=\frac{\ell_3(G)-6\ell_2(G)+16m}{6}.
\tag{7.5}
\]
Then
\[
\begin{aligned}
M_3^{(2)}(G)={}&
 c_2\bigl(m+3m(m-1)\bigr)
 +6c_{222}\binom m3 \\
&+6(n-9)p(G)(m-2)-6(n-11)a(G)
 +12\bigl(r(G)+7t(G)\bigr).
\end{aligned}
\tag{7.6}
\]
Consequently the ordinary Laplacian spectrum together with the cubic
$(n-2,2)$-moment determines the integer
\[
r(G)+7t(G).
\tag{7.7}
\]
\end{proposition}

\begin{proof}
First,
\[
\tr(L_G^3)=\sum_v d_v^3+3\sum_vd_v^2-6t(G),
\]
while
\[
s(G)=\sum_v\binom{d_v}{3}
=\frac{1}{6}\left(\sum_vd_v^3-3\sum_vd_v^2+4m\right).
\]
Together with $\ell_2=\sum_vd_v^2+2m$, this gives (7.5).

Next count pairs consisting of an adjacent unordered pair of edges and a
third distinct edge.  A triangle and a claw each contribute $3$, a
three-edge path contributes $2$, a copy of $P_3\sqcup K_2$ contributes $1$,
and a three-matching contributes $0$.  Hence
\[
p(G)(m-2)=3t(G)+3s(G)+2r(G)+q(G).
\tag{7.8}
\]
Use (7.8) to eliminate $q(G)$ from (7.4), and then use
$d_3=\binom m3-t-s-r-q$ as already incorporated there.  Since
\[
c_4-c_{222}=2n-16,\qquad
c_{32}-c_{222}=n-9,
\]
the terms depending on $s,r,t$ simplify to
\[
-6(n-11)s+12r+6(n+3)t.
\]
Writing $s=a+t$ turns this into
\[
-6(n-11)a+12(r+7t),
\]
which proves (7.6).  All quantities in (7.6) except $r+7t$ are determined
by the Laplacian spectrum and $M_3^{(2)}(G)$, proving (7.7).
\end{proof}

\end{document}
